\uuid{7bLK}
\titre{Théorème de Rolle et des accroissements finis}
\auteur{bodin}
\organisation{exo7}
\licence{CC BY-NC-SA 4.0}
\datecreate{2020-01-21}
\niveau{L1}
\module{Analyse}
\chapitre{Dérivabilité des fonctions réelles}
\sousChapitre{Théorème de Rolle et accroissements finis}
\difficulte{}
\theme{Réels}
\exo7qcmid{arnaud.arnaud.028}

\contenu{
\begin{question}
\qclassification{124.02}
Quelles sont les assertions vraies ?
\begin{answers}
    \bad{Si $f$ est dérivable sur $[a,b]$ avec $f(a) = f(b)$ il existe un unique $c \in ]a,b[$ tel que $f'( c ) = 0$.}
    \good{Si $f$ est une fonction continue sur $[a,b]$ et dérivable sur $]a,b[$ et $f'(x)$ tend vers $\ell$ quand  $x$ tend vers $a$
alors $f$ est dérivable en $a$ et $f'(a) = \ell$.}
    \bad{Soit $f(x) = \ln x$ si $x>0$ et $f(0)= 0$. Pour $x>0$ il existe $c\in]0,x[$ tel que $\ln x = \frac xc$.}
    \good{Si $f$ est dérivable sur $\Rr$ et $\lim f(x) = +1$ quand $x\rightarrow +\infty$ et $\lim f(x) = +1$ quand $x\rightarrow -\infty$ alors il existe $c \in \Rr$ tel que $f'(c)=0$.}
    \good{$\forall x > 0\ e^x \le xe^x+1$.}
\end{answers}
\end{question}
}
